\documentclass[11pt]{article}

% -------------------------------------------------
% Packages
% -------------------------------------------------
\usepackage[margin=1in]{geometry}
\usepackage{amsmath,amssymb,amsthm}
\usepackage{mathtools}
\usepackage{bm}
\usepackage{enumitem}
\usepackage[
    colorlinks=true,
    linkcolor=blue,
    urlcolor=blue,
    citecolor=blue
]{hyperref}
\usepackage{xcolor}
\usepackage{verbatim}

% -------------------------------------------------
% Formatting
% -------------------------------------------------
\setlength{\parindent}{0pt}
\setlength{\parskip}{0.3cm}

% -------------------------------------------------
% Theorem Environments
% -------------------------------------------------
\newtheorem{theorem}{Theorem}[section]
\newtheorem{definition}[theorem]{Definition}
\newtheorem{proposition}[theorem]{Proposition}
\newtheorem{remark}[theorem]{Remark}
\newtheorem{example}[theorem]{Example}

% -------------------------------------------------
% Title
% -------------------------------------------------
\title{
\textbf{Diagonalization of Defective Matrices via the Jordan Canonical Form}\\[0.3cm]
\large Theoretical Development and Practical Intuition
}

\author{Biajid Chowdhury}
\date{\today}

\begin{document}

\maketitle

% =================================================
% ABSTRACT
% =================================================

\begin{abstract}

Diagonalizing a square matrix is one of the canonical procedures in
computational mathematics.  When a matrix possesses a complete set of
linearly independent eigenvectors, a suitable change of basis transforms
the matrix into a diagonal matrix.  In that basis, the action of a
complicated linear operator is reduced to independent scalar
multiplications.  This simplification has important consequences in
linear algebra, differential equations, dynamical systems, numerical
computation, and many other areas of scientific studies.

However, not every matrix possesses enough linearly independent
eigenvectors to be diagonalized.  Such matrices are called
\emph{defective}.  At first sight, this appears to bring the
diagonalization procedure to an abrupt end: the characteristic
polynomial may provide all the eigenvalues required by the dimension of
the space, while the corresponding eigenspaces fail to provide enough
vectors to construct an eigenvector basis.

When a familiar mathematical procedure reaches such an obstruction,
a new idea is often required.  The Jordan canonical form provides one
of the most elegant examples of such an idea.  Instead of abandoning
the simplification offered by eigenvalues and eigenvectors, the theory
extends the notion of an eigenvector to that of a
\emph{generalized eigenvector}.  These vectors organize themselves into
chains that reveal the internal structure hidden behind repeated
eigenvalues.  In an appropriate basis, a defective matrix can then be
transformed into a matrix that is not necessarily diagonal, but is as
close to diagonal as its algebraic structure permits.

The purpose of this note is to develop the Jordan canonical form from
this point of view.  Rather than presenting Jordan form merely as an
algorithm for arranging eigenvalues and ones into blocks, we seek to
understand why generalized eigenvectors arise, why Jordan chains have
their particular structure, and why the resulting similarity
transformation works.  Particular emphasis is placed on practical
intuition so that the underlying ideas remain accessible even when the
algebra becomes technically demanding.

\end{abstract}

\tableofcontents
\newpage

% =================================================
\section{Introduction}
% =================================================

For a diagonalizable matrix \(A\in\mathbb{F}^{n\times n}\), there exists
an invertible matrix \(P\) such that

\[
P^{-1}AP=D,
\]

where \(D\) is diagonal.  If

\[
P=
\begin{pmatrix}
\vert & & \vert\\
v_1 & \cdots & v_n\\
\vert & & \vert
\end{pmatrix},
\]

then the columns \(v_1,\ldots,v_n\) form a basis of eigenvectors of
\(A\), and

\[
Av_i=\lambda_i v_i.
\]

Consequently,

\[
D=
\operatorname{diag}
(\lambda_1,\ldots,\lambda_n).
\]

The importance of this representation is difficult to overstate.
A matrix that may originally mix many coordinates simultaneously is
replaced, in the eigenvector basis, by independent scalar actions.
For example,

\[
D^k=
\operatorname{diag}
(\lambda_1^k,\ldots,\lambda_n^k),
\]

and

\[
e^{tD}
=
\operatorname{diag}
(e^{\lambda_1t},\ldots,e^{\lambda_nt}).
\]

Thus powers of matrices, matrix exponentials, systems of differential
equations, and many questions concerning long-term dynamical behavior
become considerably easier once a suitable eigenvector basis is known.

There is, however, a fundamental difficulty.

Suppose an \(n\times n\) matrix has \(n\) eigenvalues when algebraic
multiplicity is counted, but the corresponding eigenspaces fail to
provide \(n\) linearly independent eigenvectors.  We can no longer
construct the matrix \(P\) entirely from ordinary eigenvectors.
Diagonalization therefore fails.

This raises the central question of this note:

\[
\boxed{
\text{What should we do when the eigenvectors run out?}
}
\]

The Jordan canonical form provides an answer by asking us not to
discard the available eigenvectors, but instead to investigate the
structure surrounding them.

Suppose now that an eigenvalue \(\lambda\) has algebraic multiplicity
\(m>1\).  Ideally, we would like to find \(m\) linearly independent
eigenvectors associated with \(\lambda\).  If this happens, the repeated
eigenvalue itself causes no difficulty for diagonalization.

The difficulty arises when the eigenspace

\[
E_\lambda=\ker(A-\lambda I)
\]

has dimension strictly smaller than the algebraic multiplicity of
\(\lambda\).  In other words,

\[
\dim\ker(A-\lambda I)<m.
\]

Then there are not enough ordinary eigenvectors associated with
\(\lambda\) to account for all \(m\) occurrences of that eigenvalue in
the characteristic polynomial.

For example, suppose that \(\lambda\) has algebraic multiplicity \(3\),
but

\[
\dim\ker(A-\lambda I)=1.
\]

Then we can find only one linearly independent ordinary eigenvector,
say \(v_1\), satisfying

\[
(A-\lambda I)v_1=0.
\]

At this point ordinary diagonalization has run out of eigenvectors.
However, rather than abandoning the eigenvalue \(\lambda\), we ask
whether there are other vectors that are connected to its eigenspace
through the operator \(A-\lambda I\).

In particular, we search for a vector \(v_2\) satisfying

\[
(A-\lambda I)v_2=v_1.
\]

The vector \(v_2\) is not an ordinary eigenvector, since
\((A-\lambda I)v_2\neq 0\).  Nevertheless, applying
\(A-\lambda I\) once sends \(v_2\) directly into the eigenspace.

If another vector is still required, we may search for \(v_3\)
satisfying

\[
(A-\lambda I)v_3=v_2.
\]

The resulting vectors form the chain

\[
v_3
\;\xrightarrow{\,A-\lambda I\,}\;
v_2
\;\xrightarrow{\,A-\lambda I\,}\;
v_1
\;\xrightarrow{\,A-\lambda I\,}\;
0.
\]

Thus,

\[
(A-\lambda I)v_1=0,
\]

\[
(A-\lambda I)v_2=v_1,
\]

and

\[
(A-\lambda I)v_3=v_2.
\]

The vector \(v_1\) is an ordinary eigenvector, while \(v_2\) and
\(v_3\) are generalized eigenvectors.  Together they form what is
called a \emph{Jordan chain} associated with the eigenvalue
\(\lambda\).

This construction provides the missing basis vectors that ordinary
eigenvectors could not supply.  As we shall see, the structure of this
chain is precisely what produces a Jordan block.
The purpose of this note is not merely to state this construction.
We will develop it gradually from the failure of ordinary
diagonalization.  We will examine the role of the operator

\[
N=A-\lambda I,
\]

study the nested null spaces

\[
\ker(N)
\subseteq
\ker(N^2)
\subseteq
\ker(N^3)
\subseteq\cdots,
\]

and show how these spaces reveal the lengths and number of Jordan
chains.  This will naturally lead to generalized eigenvectors, Jordan
blocks, the minimal polynomial, and finally the Jordan canonical form.

Throughout the development, theoretical arguments will be accompanied
by practical interpretations and explicit examples.  The objective is
to make the Jordan construction feel less like a collection of
algebraic rules and more like a natural response to the failure of
ordinary diagonalization.  Even when the mathematics becomes
technically involved, the reader should be able to retain a clear
picture of what the matrix is doing and why each new construction is
needed.

% =================================================
\section{Prerequisites}
% =================================================

This note assumes that the reader has already mastered the basic
theory of eigenvalues, eigenvectors, and diagonalization.  In
particular, the reader should be comfortable with the eigenvalue
equation

\[
Av=\lambda v,
\]

and its equivalent homogeneous system

\[
(A-\lambda I)v=0.
\]

The reader should also know how to obtain eigenvalues from the
characteristic equation

\[
\det(A-\lambda I)=0,
\]

and how to determine the eigenspace associated with an eigenvalue
\(\lambda\),

\[
E_\lambda
=
\ker(A-\lambda I).
\]

We further assume familiarity with algebraic and geometric
multiplicity.  If an eigenvalue \(\lambda\) occurs \(m\) times as a
root of the characteristic polynomial, then its algebraic multiplicity
is \(m\).  Its geometric multiplicity is

\[
\dim\ker(A-\lambda I).
\]

For every eigenvalue,

\[
1
\leq
\text{geometric multiplicity}
\leq
\text{algebraic multiplicity}.
\]

The reader should also be familiar with ordinary diagonalization.  If
an \(n\times n\) matrix \(A\) possesses \(n\) linearly independent
eigenvectors \(v_1,\ldots,v_n\), then

\[
P=
\begin{pmatrix}
\vert & & \vert\\
v_1 & \cdots & v_n\\
\vert & & \vert
\end{pmatrix}
\]

is invertible and

\[
P^{-1}AP=D,
\]

where \(D\) is diagonal.

Equivalently,

\[
AP=PD.
\]

This latter form is particularly useful because it exposes what
diagonalization actually means:

\[
A
\begin{pmatrix}
\vert & & \vert\\
v_1 & \cdots & v_n\\
\vert & & \vert
\end{pmatrix}
=
\begin{pmatrix}
\vert & & \vert\\
v_1 & \cdots & v_n\\
\vert & & \vert
\end{pmatrix}
\begin{pmatrix}
\lambda_1 & & 0\\
& \ddots &\\
0 & & \lambda_n
\end{pmatrix}.
\]

Thus diagonalization is fundamentally a search for a basis in which
the action of the linear operator becomes as simple as possible.

Finally, basic familiarity with null spaces, rank, basis, linear
independence, similarity transformations, and matrix multiplication is
assumed.

No previous knowledge of generalized eigenvectors, Jordan chains,
minimal polynomials, or the Jordan canonical form is required.  These
ideas will be developed from the beginning.


% =================================================
\section{Historical Development of the Jordan Canonical Form}
% =================================================

The Jordan canonical form did not emerge as an isolated construction.
It grew out of a much broader nineteenth-century effort to understand
linear transformations through their characteristic quantities and to
find coordinate systems in which those transformations take their
simplest possible forms.

\subsection{The Road Toward Eigenvalue Decomposition}

The ideas that we now organize under the language of eigenvalues,
eigenvectors, and matrix diagonalization developed gradually over a
long period of time.  Their origins can be traced to problems involving
systems of linear equations, differential equations, quadratic forms,
mechanics, and astronomy.

During the eighteenth century, mathematicians such as
\textbf{Leonhard Euler} and \textbf{Joseph-Louis Lagrange} encountered
what we would now recognize as eigenvalue problems while studying
systems of differential equations and small oscillations in mechanical
systems.  In such problems, certain special modes of motion allowed a
complicated coupled system to be separated into simpler independent
motions.

The subject advanced considerably during the nineteenth century.
\textbf{Augustin-Louis Cauchy} developed important results involving
determinants, characteristic equations, and quadratic forms.  These
ideas helped establish the connection between characteristic roots and
the simplification of linear transformations.

At the same time, the theory of determinants and linear substitutions
was becoming increasingly systematic.  Mathematicians were gradually
recognizing a central principle: although the numerical entries used
to represent a linear transformation depend on the chosen coordinates,
certain underlying properties of the transformation remain unchanged.

This naturally led to the search for special coordinate systems in
which a linear transformation could be represented as simply as
possible.

In modern language, one of the most desirable outcomes of this search
is \emph{eigenvalue decomposition}, or diagonalization.  When enough
independent eigenvectors exist, the transformation can be represented
by a diagonal matrix.  A complicated coupled transformation is thereby
reduced to a collection of independent scalar actions.

This was an extraordinarily powerful idea.  But it came with an
important limitation: not every linear transformation possesses enough
independent eigenvectors to permit such a complete diagonal reduction.

The question therefore became broader than diagonalization itself:

\begin{center}
\emph{If complete diagonal reduction is impossible, what is the
simplest form to which a linear transformation can still be reduced?}
\end{center}

It was in the development of this broader problem that Camille Jordan
entered the story.

\subsection{Camille Jordan and Canonical Reduction}

\textbf{Marie Ennemond Camille Jordan} (1838--1922) was a French
mathematician whose work had a major influence on nineteenth-century
algebra.  His contributions extended well beyond the subject discussed
in this note and included important work in group theory, algebra,
analysis, and geometry.

Jordan was particularly interested in the theory of substitutions.
In the terminology of the period, a linear substitution described what
we would now naturally regard as a linear transformation represented
through a system of linear expressions.

The fundamental problem was one of classification and reduction.
Rather than accepting the potentially complicated original
representation of a transformation, mathematicians wanted to determine
whether it could be converted into a simpler standard form without
losing its essential algebraic information.

Jordan developed a systematic theory for carrying this reduction much
further, including precisely those situations in which a complete
diagonal reduction was unavailable.

His treatment appeared prominently in his monumental 1870 work

\begin{center}
\emph{Trait\'e des substitutions et des \'equations alg\'ebriques}.
\end{center}

The canonical structure associated with this work eventually became
known as the \emph{Jordan canonical form}, or \emph{Jordan normal
form}.

Jordan's contribution was important because it showed that the failure
of diagonalization was not the end of the simplification process.
Even when a transformation could not be completely decomposed into
independent eigendirections, its remaining algebraic structure could
still be organized into a remarkably simple and systematic canonical
form.

\subsection{Jordan's Place in a Larger Mathematical Development}

As is often the case with major mathematical theories, the Jordan
canonical form should not be understood as the product of a single
mathematician working in isolation.  It emerged from a long development
involving characteristic equations, determinants, linear
transformations, invariant quantities, and canonical reduction.

Important related work was carried out by mathematicians including
Cauchy and, independently in the theory of elementary divisors,
\textbf{Karl Weierstrass}.  The nineteenth century was a period in
which many of the ideas that would eventually become modern linear
algebra were being developed simultaneously from several directions.

Jordan's achievement was to provide a particularly systematic
canonical theory for linear substitutions.  The language and notation
used today came later, as matrix theory and abstract linear algebra
were organized into their modern forms.

Consequently, it would be historically misleading to imagine Jordan
sitting down with the modern formula

\[
A=PDP^{-1}
\]

and attempting to repair it when it failed.  That is a useful modern
way for us to motivate the subject, but it is not a literal description
of its historical development.

Historically, the movement was broader:

\[
\boxed{
\begin{gathered}
\text{characteristic equations and special modes}\\
\longrightarrow\\
\text{eigenvalue-type ideas and linear substitutions}\\
\longrightarrow\\
\text{search for simpler representations}\\
\longrightarrow\\
\text{canonical reduction of transformations}\\
\longrightarrow\\
\text{Jordan's canonical theory}.
\end{gathered}
}
\]

Today, this history appears in a particularly elegant form in a first
course in linear algebra.  We first learn that eigenvectors can provide
a basis in which a matrix becomes diagonal.  We then discover that some
matrices do not possess enough eigenvectors.  Jordan canonical form
shows that the nineteenth-century search for simplification does not
stop there.

With this historical perspective in place, we now return to the
mathematics.  The first theoretical tool we shall need is the
Cayley--Hamilton theorem.
% =================================================


% =================================================
\section{Background: The Cayley--Hamilton Theorem}
% =================================================

Before developing the theory of generalized eigenvectors and Jordan
chains, we recall one result that will play an important role throughout
the discussion: the \emph{Cayley--Hamilton theorem}.

The theorem establishes a remarkable connection between a matrix and
its own characteristic polynomial.

\subsection{The Characteristic Polynomial}

Let \(A\in\mathbb{F}^{n\times n}\).  We define the characteristic
polynomial of \(A\) by

\[
p_A(\lambda)=\det(\lambda I-A).
\]

Since \(A\) is an \(n\times n\) matrix, \(p_A(\lambda)\) is a polynomial
of degree \(n\).  It can therefore be written in the form

\[
p_A(\lambda)
=
\lambda^n
+c_{n-1}\lambda^{n-1}
+\cdots
+c_1\lambda
+c_0.
\]

The eigenvalues of \(A\) are precisely the roots of this polynomial.

For example, suppose that a \(3\times3\) matrix has characteristic
polynomial

\[
p_A(\lambda)=(\lambda-2)^3.
\]

Then \(2\) is the only eigenvalue of \(A\), and it has algebraic
multiplicity \(3\).

At this stage, however, the characteristic polynomial appears to be a
statement about the scalar variable \(\lambda\).  The Cayley--Hamilton
theorem tells us something much stronger: the matrix \(A\) itself
satisfies this polynomial.

\subsection{The Cayley--Hamilton Theorem}

\begin{theorem}[Cayley--Hamilton Theorem]
Every square matrix satisfies its own characteristic polynomial.

More precisely, if

\[
p_A(\lambda)
=
\lambda^n
+c_{n-1}\lambda^{n-1}
+\cdots
+c_1\lambda
+c_0,
\]

then

\[
\boxed{
p_A(A)
=
A^n
+c_{n-1}A^{n-1}
+\cdots
+c_1A
+c_0I
=
0.
}
\]

Here \(0\) denotes the zero matrix.
\end{theorem}

The notation \(p_A(A)\) simply means that the scalar variable
\(\lambda\) in the characteristic polynomial is replaced by the
matrix \(A\).  The constant term must consequently be interpreted as
a multiple of the identity matrix.

For example, if

\[
p_A(\lambda)
=
\lambda^3-6\lambda^2+12\lambda-8,
\]

then the Cayley--Hamilton theorem gives

\[
A^3-6A^2+12A-8I=0.
\]

Thus the matrix obeys the same algebraic relation as the roots of its
characteristic polynomial.

\subsection{A Particularly Important Repeated-Eigenvalue Case}

The connection with defective matrices becomes especially interesting
when the characteristic polynomial contains a repeated eigenvalue.

Suppose, for example, that

\[
p_A(\lambda)=(\lambda-2)^3.
\]

The Cayley--Hamilton theorem immediately gives

\[
p_A(A)=0,
\]

and hence

\[
\boxed{
(A-2I)^3=0.
}
\]

More generally, if an \(n\times n\) matrix has only one eigenvalue
\(\lambda_0\), with algebraic multiplicity \(n\), then

\[
p_A(\lambda)
=
(\lambda-\lambda_0)^n,
\]

and therefore

\[
\boxed{
(A-\lambda_0 I)^n=0.
}
\]

This observation suggests introducing the shifted matrix

\[
N=A-\lambda_0 I.
\]

Then

\[
A=\lambda_0 I+N
\]

and Cayley--Hamilton tells us that

\[
N^n=0.
\]

A matrix whose sufficiently high power is zero is called
\emph{nilpotent}.  Thus, in this special situation, subtracting the
eigenvalue from the matrix exposes a nilpotent part:

\[
\boxed{
A=\lambda_0 I+N,
\qquad
N^n=0.
}
\]

This decomposition will become extremely important when we develop the
Jordan canonical form.

\subsection{What Does \(N^n=0\) Actually Mean?}

The matrix equation

\[
N^n=0
\]

has a useful interpretation in terms of vectors.

For every vector \(v\),

\[
N^n v=0.
\]

Thus repeated application of \(N\) must eventually send every vector
to the zero vector.

For example, if

\[
N^3=0,
\]

then for every vector \(v\),

\[
N^3v=0.
\]

This does \emph{not} mean that every vector requires exactly three
applications of \(N\) before reaching zero.

Some vectors may already satisfy

\[
Nv=0.
\]

Others may satisfy

\[
Nv\neq0,
\qquad
N^2v=0,
\]

while still others may satisfy

\[
Nv\neq0,
\qquad
N^2v\neq0,
\qquad
N^3v=0.
\]

Therefore, the statement

\[
N^3=0
\]

should be interpreted as saying that three applications are
\emph{sufficient} for every vector, not that three applications are
\emph{necessary} for every vector.

Schematically, different vectors may behave as

\[
v_1
\xrightarrow{\,N\,}
0,
\]

or

\[
v_2
\xrightarrow{\,N\,}
v_1
\xrightarrow{\,N\,}
0,
\]

or

\[
v_3
\xrightarrow{\,N\,}
v_2
\xrightarrow{\,N\,}
v_1
\xrightarrow{\,N\,}
0.
\]

This simple observation will later become one of the central ideas
behind generalized eigenvectors and Jordan chains.

\subsection{Why Cayley--Hamilton Matters for Our Development}

The Cayley--Hamilton theorem does more than provide a convenient
identity involving powers of a matrix.  It tells us that the
characteristic polynomial places an algebraic restriction on the
action of the matrix itself.

In the repeated-eigenvalue example

\[
p_A(\lambda)=(\lambda-2)^3,
\]

we know that

\[
(A-2I)^3=0.
\]

At the same time, an ordinary eigenvector \(v_1\) associated with
\(\lambda=2\) satisfies

\[
(A-2I)v_1=0.
\]

These two statements naturally raise an important question:

\[
\boxed{
\begin{gathered}
(A-2I)^3 \text{ annihilates every vector,}\\
\text{but an eigenvector is annihilated after only one application.}\\[2mm]
\text{What's the logic hiding here?}
\end{gathered}
}
\]

Investigating precisely these intermediate possibilities will lead us
to generalized eigenvectors, nested null spaces, Jordan chains, and
ultimately the Jordan canonical form.


% =================================================
% =================================================
\section{From Ordinary Eigenvectors to Generalized Eigenvectors}
% =================================================

We are now ready to begin the main development of the Jordan canonical
form.

Suppose that \(\lambda\) is an eigenvalue of \(A\), but its eigenspace
does not contain enough linearly independent eigenvectors to construct
a basis.  Ordinary diagonalization then fails.

The question is not whether the missing directions disappear, but
whether they can be recovered in a more general form.

Define

\[
\boxed{
N=A-\lambda I.
}
\]

An ordinary eigenvector \(v\neq 0\) associated with \(\lambda\)
satisfies

\[
Nv=0.
\]

Hence the ordinary eigenspace is

\[
\boxed{
E_\lambda
=
\ker(A-\lambda I)
=
\ker N.
}
\]

Thus ordinary eigenvectors are precisely the nonzero vectors that are
annihilated after one application of \(N\).


% -------------------------------------------------
\subsection{Looking Beyond the Ordinary Eigenspace}
% -------------------------------------------------

To see where additional vectors may come from, consider first the
special case in which the characteristic polynomial is

\[
p_A(x)=(x-\lambda)^m.
\]

By the Cayley--Hamilton theorem,

\[
(A-\lambda I)^m=0,
\]

and therefore

\[
\boxed{
N^m=0.
}
\]

Consequently,

\[
N^m v=0
\]

for every vector \(v\).

This is a much weaker condition than

\[
Nv=0.
\]

An ordinary eigenvector is annihilated immediately, but other vectors
may require several applications of \(N\) before reaching zero.  This
suggests studying the sequence

\[
\boxed{
\ker N
\subseteq
\ker N^2
\subseteq
\ker N^3
\subseteq
\cdots
\subseteq
\ker N^m.
}
\]

The reason these spaces are nested is simple.  If

\[
v\in\ker N^j,
\]

then

\[
N^jv=0.
\]

Applying \(N\) once more gives

\[
N^{j+1}v=N(N^jv)=0,
\]

and hence

\[
v\in\ker N^{j+1}.
\]

Therefore, for every \(j\geq1\),

\[
\boxed{
\ker N^j\subseteq\ker N^{j+1}.
}
\]

These nested null spaces reveal structure that cannot be seen from the
ordinary eigenspace alone.

\begin{remark}[Do Not Reverse the Inclusion]
The implication

\[
N^jv=0
\quad\Longrightarrow\quad
N^{j+1}v=0
\]

always holds, but the converse need not hold.

A vector may satisfy

\[
N^{j+1}v=0
\]

while

\[
N^jv\neq0.
\]

Thus

\[
\ker N^j
\subseteq
\ker N^{j+1}
\]

does not imply equality.  The higher null spaces may contain genuinely
new vectors.

The direction to remember is therefore

\[
\boxed{
Nv=0
\Longrightarrow
N^2v=0
\Longrightarrow
N^3v=0
\Longrightarrow
\cdots,
}
\]

and not the other way around.
\end{remark}


% -------------------------------------------------
\subsection{Generalized Eigenvectors}
% -------------------------------------------------

The nested null spaces motivate an extension of the ordinary
eigenvector concept.

\begin{definition}[Generalized Eigenvector]
Let \(\lambda\) be an eigenvalue of \(A\).  A nonzero vector \(v\) is
called a \emph{generalized eigenvector} associated with \(\lambda\) if

\[
\boxed{
(A-\lambda I)^k v=0
}
\]

for some positive integer \(k\).

Equivalently,

\[
N^k v=0
\]

for some \(k\geq1\).
\end{definition}

Ordinary eigenvectors are therefore included automatically: they are
the special case \(k=1\).

More precisely, if \(k\) is the smallest positive integer for which

\[
N^k v=0,
\]

then

\[
N^{k-1}v\neq0.
\]

Such a vector is called a generalized eigenvector of \emph{rank \(k\)}.
Thus

\[
\boxed{
N^k v=0,
\qquad
N^{k-1}v\neq0
}
\]

means that \(v\) requires exactly \(k\) applications of \(N\) to reach
zero.

Equivalently,

\[
\boxed{
v\in\ker N^k\setminus\ker N^{k-1}.
}
\]

This gives a natural hierarchy:

\[
\begin{array}{ccl}
\ker N
&:&
\text{ordinary eigenvectors},\\[1mm]
\ker N^2\setminus\ker N
&:&
\text{rank-2 generalized eigenvectors},\\[1mm]
\ker N^3\setminus\ker N^2
&:&
\text{rank-3 generalized eigenvectors},\\
&\vdots&\\
\ker N^k\setminus\ker N^{k-1}
&:&
\text{rank-\(k\) generalized eigenvectors}.
\end{array}
\]


% -------------------------------------------------
\subsection{From Generalized Eigenvectors to a Chain}
% -------------------------------------------------

Generalized eigenvectors become especially useful when they are linked
by successive applications of \(N\).

Suppose \(v_k\) is a generalized eigenvector of rank \(k\).  Define

\[
v_{k-1}=Nv_k,
\]

then

\[
v_{k-2}=Nv_{k-1}=N^2v_k,
\]

and continue in this way until

\[
v_1=N^{k-1}v_k.
\]

Because \(v_k\) has rank \(k\),

\[
v_1\neq0,
\]

while

\[
Nv_1=N^kv_k=0.
\]

Hence \(v_1\) is an ordinary eigenvector.

The resulting vectors satisfy

\[
\boxed{
Nv_1=0,
\qquad
Nv_j=v_{j-1},
\quad
j=2,\ldots,k.
}
\]

Schematically,

\[
\boxed{
v_k
\xrightarrow{\,N\,}
v_{k-1}
\xrightarrow{\,N\,}
\cdots
\xrightarrow{\,N\,}
v_2
\xrightarrow{\,N\,}
v_1
\xrightarrow{\,N\,}
0.
}
\]

Thus a rank-\(k\) generalized eigenvector naturally generates a
sequence of vectors descending through the nested null spaces.


% -------------------------------------------------
\subsection{Jordan Chains}
% -------------------------------------------------

The preceding sequence is the central structure behind the Jordan
canonical form.

\begin{definition}[Jordan Chain]
Let \(\lambda\) be an eigenvalue of \(A\).  A sequence of nonzero
vectors

\[
v_1,v_2,\ldots,v_k
\]

is called a \emph{Jordan chain of length \(k\)} associated with
\(\lambda\) if

\[
\boxed{
(A-\lambda I)v_1=0,
\qquad
(A-\lambda I)v_j=v_{j-1},
\quad
j=2,\ldots,k.
}
\]

Equivalently, with \(N=A-\lambda I\),

\[
Nv_1=0,
\qquad
Nv_j=v_{j-1},
\quad
j=2,\ldots,k.
\]
\end{definition}

The chain therefore has the form

\[
\boxed{
v_k
\xrightarrow{\,N\,}
v_{k-1}
\xrightarrow{\,N\,}
\cdots
\xrightarrow{\,N\,}
v_2
\xrightarrow{\,N\,}
v_1
\xrightarrow{\,N\,}
0.
}
\]

The vector \(v_1\) is an ordinary eigenvector, while \(v_j\) is a
generalized eigenvector of rank \(j\).  Indeed,

\[
N^jv_j=0,
\qquad
N^{j-1}v_j=v_1\neq0.
\]

Therefore,

\[
\boxed{
v_j\in\ker N^j\setminus\ker N^{j-1}.
}
\]

A Jordan chain can thus be viewed as a sequence that connects the
successive levels of the nested null spaces.


% -------------------------------------------------
\subsection{Why a Jordan Chain Produces a Jordan Block}
% -------------------------------------------------

The defining relations of a Jordan chain immediately determine how
\(A\) acts on its vectors.

Since

\[
N=A-\lambda I,
\]

The first chain relation,

\[
Nv_1=0,
\]

gives, after substituting

\[
N=A-\lambda I,
\]

the equation

\[
(A-\lambda I)v_1=0.
\]

Expanding,

\[
Av_1-\lambda Iv_1=0.
\]

Since

\[
Iv_1=v_1,
\]

we obtain

\[
Av_1-\lambda v_1=0,
\]

and therefore

\[
\boxed{
Av_1=\lambda v_1.
}
\]

while

\[
Nv_j=v_{j-1},
\qquad j=2,\ldots,k,
\]

gives, after substituting

\[
N=A-\lambda I,
\]

the relation

\[
(A-\lambda I)v_j=v_{j-1}.
\]

Expanding the left-hand side,

\[
Av_j-\lambda Iv_j=v_{j-1}.
\]

Since

\[
Iv_j=v_j,
\]

we obtain

\[
Av_j-\lambda v_j=v_{j-1}.
\]

Therefore,

\[
\boxed{
Av_j=\lambda v_j+v_{j-1},
\qquad j=2,\ldots,k.
}
\]

Thus the action of \(A\) on the entire Jordan chain is

\[
\boxed{
\begin{aligned}
Av_1 &= \lambda v_1,\\
Av_2 &= \lambda v_2+v_1,\\
Av_3 &= \lambda v_3+v_2,\\
&\ \vdots\\
Av_j &= \lambda v_j+v_{j-1},\\
&\ \vdots\\
Av_k &= \lambda v_k+v_{k-1}.
\end{aligned}
}
\]
Relative to the ordered chain

\[
\mathcal{B}=(v_1,v_2,\ldots,v_k),
\]

these relations are represented by the matrix

\[
\boxed{
J_k(\lambda)
=
\begin{pmatrix}
\lambda & 1 & 0 & \cdots & 0\\
0 & \lambda & 1 & \ddots & \vdots\\
\vdots & \ddots & \ddots & \ddots & 0\\
0 & \cdots & 0 & \lambda & 1\\
0 & \cdots & \cdots & 0 & \lambda
\end{pmatrix}.
}
\]

This matrix is called the \emph{Jordan block of size \(k\)}
associated with the eigenvalue \(\lambda\).

If the chain vectors form a basis of the space under consideration,
let

\[
P=
\begin{pmatrix}
\vert & \vert & & \vert\\
v_1 & v_2 & \cdots & v_k\\
\vert & \vert & & \vert
\end{pmatrix}.
\]

Then the chain relations give

\[
AP=PJ_k(\lambda),
\]

and therefore, provided \(P\) is invertible,

\[
\boxed{
P^{-1}AP=J_k(\lambda).
}
\]

The Jordan block is therefore not an arbitrary substitute for a
diagonal matrix.  It is the matrix representation naturally produced
when \(A\) is expressed in a basis organized as a Jordan chain.


% -------------------------------------------------
\subsection{Where Do the Ones Come From?}
% -------------------------------------------------

The superdiagonal \(1\)'s in a Jordan block are a direct consequence
of the chain relation

\[
Av_j=\lambda v_j+v_{j-1}.
\]

Relative to the ordered basis

\[
\mathcal{B}=(v_1,v_2,\ldots,v_k),
\]

the vector \(Av_j\) therefore has coefficient \(\lambda\) in the
\(v_j\)-direction and coefficient \(1\) in the \(v_{j-1}\)-direction.

Consequently,

\[
\boxed{
\lambda\text{ on the diagonal}
\quad\longleftrightarrow\quad
\text{the eigenvalue contribution},
}
\]

whereas

\[
\boxed{
1\text{ on the superdiagonal}
\quad\longleftrightarrow\quad
\text{the connection }v_j\mapsto v_{j-1}.
}
\]

Thus the \(1\)'s are not inserted merely by convention.  They record
the defining structure of the Jordan chain:

\[
\boxed{
(A-\lambda I)v_j=v_{j-1}.
}
\]

This is the fundamental reason that a Jordan block has its
characteristic form.

If the reader is still having difficulty seeing where the \(1\)'s come
from, recall the basic rule of matrix--vector multiplication.  The
\(1\) selects the preceding vector \(v_{j-1}\), while \(\lambda\)
selects the current vector \(v_j\), giving precisely

\[
Av_j=v_{j-1}+\lambda v_j.
\]

If this is still not evident, we recommend reviewing our other notes,
where matrix--vector and matrix--matrix multiplication are discussed
explicitly from several different perspectives.

% =================================================
% =================================================
\section{Practical Intuition Behind Jordan Chains}
% =================================================

So far, we have developed the mathematics behind generalized
eigenvectors, Jordan chains, and Jordan blocks in considerable detail.
We have used the Cayley--Hamilton theorem, nested null spaces, and the
operator \(A-\lambda I\) to explain how the Jordan structure arises.

We now temporarily put most of that formal mathematics aside.

The purpose of this section is to develop an intuitive picture of the
same phenomenon through familiar situations.  If the reader found the
previous theoretical development difficult, it may actually be useful
to begin here, develop the intuition, and then return to the preceding
section.  The algebra may look considerably more natural the second
time around.

The central puzzle that might remain in our mind comes from comparing
two facts.

In the single-eigenvalue situation considered earlier, the
Cayley--Hamilton theorem tells us that, for some positive integer \(k\),

\[
\boxed{
(A-\lambda I)^k v=0
}
\]

for every vector \(v\).

Yet an ordinary eigenvector \(v_1\) satisfies the much stronger
condition

\[
\boxed{
(A-\lambda I)v_1=0.
}
\]

Something interesting is therefore happening.  The operator may require
several applications before an arbitrary vector is guaranteed to
disappear, yet this particular vector disappears after only one.

Let us forget matrices for a moment and imagine an everyday situation.

% -------------------------------------------------
\subsection{The Medicine Analogy}
% -------------------------------------------------

Suppose a physician prescribes a medicine and determines that a patient
may require several doses before the treatment is complete.

For the sake of the analogy, suppose that as many as \(k\) doses may
be required.

We therefore expect something like

\[
\text{patient}
\longrightarrow
\text{dose}
\longrightarrow
\text{dose}
\longrightarrow
\cdots
\longrightarrow
\text{recovery}.
\]

But something surprising happens.

One particular patient becomes well immediately after the first dose.

Naturally, we would ask:

\begin{center}
\emph{Why did this patient require only one dose when several doses
might have been necessary?}
\end{center}

One possible explanation is that the first dose we observed was not
really the beginning of the story.  Perhaps the patient had already
received the same medicine through some other route.  It might have
been administered previously without our knowledge, or an equivalent
substance might already have entered the patient's system through some
other source.

The important point is not the medical details of the analogy.  The
important point is that an unexpectedly early outcome makes us suspect
that there may be a \emph{history behind what we are currently
observing}.

This is precisely the kind of intuition that is useful for the Jordan
construction.

The Cayley--Hamilton theorem tells us that repeated application of

\[
A-\lambda I
\]

will eventually annihilate the relevant vectors.  But the ordinary
eigenvector \(v_1\) is already annihilated after one application:

\[
(A-\lambda I)v_1=0.
\]

Instead of merely accepting this as the end of the story, we may ask
whether \(v_1\) itself could have come from somewhere.

Could there be another vector \(v_2\) such that

\[
(A-\lambda I)v_2=v_1?
\]

If so, \(v_1\) has what we may informally think of as a
\emph{parent}.  And then another question immediately arises:

\[
\text{Could }v_2\text{ itself have a parent?}
\]

Perhaps there exists \(v_3\) such that

\[
(A-\lambda I)v_3=v_2.
\]

Continuing this way, what initially appeared to be a single ordinary
eigenvector may turn out to be the visible end of a much longer
history:

\[
\boxed{
v_k
\longrightarrow
v_{k-1}
\longrightarrow
\cdots
\longrightarrow
v_2
\longrightarrow
v_1
\longrightarrow
0.
}
\]

The ordinary eigenvector \(v_1\) is therefore not necessarily the
beginning of the story.  It may be the \emph{last surviving member} of
a sequence of vectors connected through repeated applications of
\(A-\lambda I\).

There is another useful way to visualize the same idea.  Instead of
thinking about medicine, consider the inheritance of a fatal genetic disease
through a family.

For the purpose of this analogy, imagine that being sent to the null
space---that is, being mapped to zero---corresponds to dying from the
disease.  A particular genetic condition may be carried through several
generations without causing that outcome.  A grandparent may carry the
relevant genetic variant, pass it to a parent, and the parent may in
turn pass it to a child.  Then, for biological reasons that may involve
many additional factors, the disease may become severe in the child, and unfortunately they die.

The unexpected event naturally encourages the physicians to look
backward through the family history.  They may discover that what
appeared to happen suddenly in the child actually has an ancestry:

\[
\boxed{
\text{grandparent}
\longrightarrow
\text{parent}
\longrightarrow
\text{child}
\longrightarrow
\text{death}.
}
\]

The child is where the outcome became visible, but the story did not
necessarily begin with the child.  Tracing backward reveals the
generations that carried the relevant inherited factor.

We may use the same picture for our vectors.  If reaching zero is
thought of as reaching the ``death'' of the process, then

\[
v_1
\xrightarrow{\,A-\lambda I\,}
0
\]

shows us only the final visible step.  We may therefore investigate
whether \(v_1\) has a parent \(v_2\), whether \(v_2\) has a parent
\(v_3\), and so on.  Tracing this mathematical family tree backward
may reveal

\[
\boxed{
v_k
\xrightarrow{\,A-\lambda I\,}
v_{k-1}
\xrightarrow{\,A-\lambda I\,}
\cdots
\xrightarrow{\,A-\lambda I\,}
v_2
\xrightarrow{\,A-\lambda I\,}
v_1
\xrightarrow{\,A-\lambda I\,}
0.
}
\]

Of course, this is only an analogy: real genetic inheritance and
disease expression are considerably more complicated.  Its purpose
here is simply to emphasize one idea.  When we observe \(v_1\) being
sent immediately to zero, we can look backward and ask whether there
is an entire ancestry of vectors behind it.

This sequence is exactly what the mathematics later identifies as a
Jordan chain.

% =================================================
\section{A Complete \(3\times3\) Example}
% =================================================

The theoretical development of the Jordan canonical form can appear
somewhat convoluted when encountered for the first time.  Fortunately,
the computational problems themselves are usually much less difficult
than the theory might suggest.

In a classroom or examination setting, a sensible problem designed to
test the basic Jordan-form construction will ordinarily involve a small
matrix.  A \(3\times3\) example is already large enough to reveal the
essential structure: repeated eigenvalues, a shortage of ordinary
eigenvectors, generalized eigenvectors, a Jordan chain, the Jordan
block, and the similarity transformation.

There is therefore little pedagogical value in beginning with an
unnecessarily large matrix.  Instead, we will solve one carefully
chosen \(3\times3\) problem completely and step by step.  This example
will reveal the rest of the story developed in the previous sections.

Consider

\[
\boxed{
A=
\begin{pmatrix}
9&7&3\\
-9&-7&-4\\
4&4&4
\end{pmatrix}.
}
\]

Our goal is to find an invertible matrix \(P\) and a Jordan matrix
\(J\) such that

\[
\boxed{
P^{-1}AP=J.
}
\]

We proceed one step at a time.


% -------------------------------------------------
\subsection{Step 1: Find the Characteristic Polynomial}
% -------------------------------------------------

We begin with

\[
p_A(x)=\det(xI-A).
\]

Since

\[
xI=
\begin{pmatrix}
x&0&0\\
0&x&0\\
0&0&x
\end{pmatrix},
\]

we obtain

\[
xI-A
=
\begin{pmatrix}
x&0&0\\
0&x&0\\
0&0&x
\end{pmatrix}
-
\begin{pmatrix}
9&7&3\\
-9&-7&-4\\
4&4&4
\end{pmatrix}.
\]

Therefore,

\[
xI-A
=
\begin{pmatrix}
x-9&-7&-3\\
9&x+7&4\\
-4&-4&x-4
\end{pmatrix}.
\]

Hence

\[
p_A(x)
=
\det
\begin{pmatrix}
x-9&-7&-3\\
9&x+7&4\\
-4&-4&x-4
\end{pmatrix}.
\]

Expanding along the first row,

\[
\begin{aligned}
p_A(x)
&=
(x-9)
\begin{vmatrix}
x+7&4\\
-4&x-4
\end{vmatrix}
-
(-7)
\begin{vmatrix}
9&4\\
-4&x-4
\end{vmatrix} \\
&\qquad
+
(-3)
\begin{vmatrix}
9&x+7\\
-4&-4
\end{vmatrix}.
\end{aligned}
\]

Now compute each \(2\times2\) determinant separately.

First,

\[
\begin{aligned}
\begin{vmatrix}
x+7&4\\
-4&x-4
\end{vmatrix}
&=
(x+7)(x-4)-4(-4)\\
&=
(x+7)(x-4)+16\\
&=
x^2+3x-28+16\\
&=
x^2+3x-12.
\end{aligned}
\]

Second,

\[
\begin{aligned}
\begin{vmatrix}
9&4\\
-4&x-4
\end{vmatrix}
&=
9(x-4)-4(-4)\\
&=
9x-36+16\\
&=
9x-20.
\end{aligned}
\]

Third,

\[
\begin{aligned}
\begin{vmatrix}
9&x+7\\
-4&-4
\end{vmatrix}
&=
9(-4)-(x+7)(-4)\\
&=
-36+4(x+7)\\
&=
-36+4x+28\\
&=
4x-8.
\end{aligned}
\]

Substituting these results,

\[
\begin{aligned}
p_A(x)
&=
(x-9)(x^2+3x-12)
+
7(9x-20)
-
3(4x-8).
\end{aligned}
\]

Expand the first product:

\[
\begin{aligned}
(x-9)(x^2+3x-12)
&=
x(x^2+3x-12)
-
9(x^2+3x-12)\\
&=
x^3+3x^2-12x
-
9x^2-27x+108\\
&=
x^3-6x^2-39x+108.
\end{aligned}
\]

Also,

\[
7(9x-20)=63x-140,
\]

and

\[
-3(4x-8)=-12x+24.
\]

Therefore,

\[
\begin{aligned}
p_A(x)
&=
x^3-6x^2-39x+108
+63x-140
-12x+24\\
&=
x^3-6x^2+12x-8.
\end{aligned}
\]

But

\[
x^3-6x^2+12x-8=(x-2)^3.
\]

Hence

\[
\boxed{
p_A(x)=(x-2)^3.
}
\]

Therefore the matrix has only one eigenvalue,

\[
\boxed{
\lambda=2,
}
\]

with algebraic multiplicity

\[
\boxed{
3.
}
\]

Throughout this calculation, we have used \(x\) rather than
\(\lambda\) as the variable in the characteristic polynomial simply
for easier typesetting and readability.  There is no mathematical
difference between the two choices.

In this context, \(x\) and \(\lambda\) are only dummy variables used
to represent the argument of the polynomial.  Thus

\[
p_A(x)=(x-2)^3
\]

and

\[
p_A(\lambda)=(\lambda-2)^3
\]

express exactly the same characteristic polynomial.  The reader should
therefore not attach any special significance to the choice of symbol.


% -------------------------------------------------
\subsection{Step 2: Find the Ordinary Eigenvectors}
% -------------------------------------------------

We now solve

\[
(A-2I)v=0.
\]

First compute

\[
A-2I
=
\begin{pmatrix}
9&7&3\\
-9&-7&-4\\
4&4&4
\end{pmatrix}
-
\begin{pmatrix}
2&0&0\\
0&2&0\\
0&0&2
\end{pmatrix}.
\]

Thus

\[
\boxed{
N=A-2I
=
\begin{pmatrix}
7&7&3\\
-9&-9&-4\\
4&4&2
\end{pmatrix}.
}
\]

Let

\[
v=
\begin{pmatrix}
x\\
y\\
z
\end{pmatrix}.
\]

Then

\[
Nv=0
\]

becomes

\[
\begin{pmatrix}
7&7&3\\
-9&-9&-4\\
4&4&2
\end{pmatrix}
\begin{pmatrix}
x\\
y\\
z
\end{pmatrix}
=
\begin{pmatrix}
0\\
0\\
0
\end{pmatrix}.
\]

Therefore,

\[
\begin{aligned}
7x+7y+3z&=0,\\
-9x-9y-4z&=0,\\
4x+4y+2z&=0.
\end{aligned}
\]

Since \(x\) and \(y\) always appear as \(x+y\), let

\[
s=x+y.
\]

The system becomes

\[
\begin{aligned}
7s+3z&=0,\\
-9s-4z&=0,\\
4s+2z&=0.
\end{aligned}
\]

From the third equation,

\[
4s+2z=0.
\]

Hence

\[
2s+z=0,
\]

so

\[
z=-2s.
\]

Substitute this into the first equation:

\[
7s+3(-2s)=0.
\]

Thus

\[
7s-6s=0,
\]

and therefore

\[
s=0.
\]

Hence

\[
z=0.
\]

Since

\[
s=x+y=0,
\]

we have

\[
y=-x.
\]

Therefore every eigenvector has the form

\[
v=
x
\begin{pmatrix}
1\\
-1\\
0
\end{pmatrix}.
\]

We choose

\[
\boxed{
v_1=
\begin{pmatrix}
1\\
-1\\
0
\end{pmatrix}.
}
\]

Thus

\[
\boxed{
E_2
=
\ker(A-2I)
=
\operatorname{span}
\left\{
\begin{pmatrix}
1\\
-1\\
0
\end{pmatrix}
\right\}.
}
\]

Therefore,

\[
\dim E_2=1.
\]

The algebraic multiplicity is \(3\), but the geometric multiplicity is
only \(1\):

\[
\boxed{
1<3.
}
\]

Hence \(A\) is defective and cannot be diagonalized.

This is exactly the situation for which the Jordan construction is
needed.


% -------------------------------------------------
\subsection{Step 3: Search for the Parent of \(v_1\)}
% -------------------------------------------------

We have

\[
Nv_1=0.
\]

Following the intuition developed earlier, we now ask whether \(v_1\)
has a ``parent.''

That is, can we find a vector \(v_2\) satisfying

\[
\boxed{
Nv_2=v_1?
}
\]

Let

\[
v_2=
\begin{pmatrix}
a\\
b\\
c
\end{pmatrix}.
\]

Then

\[
\begin{pmatrix}
7&7&3\\
-9&-9&-4\\
4&4&2
\end{pmatrix}
\begin{pmatrix}
a\\
b\\
c
\end{pmatrix}
=
\begin{pmatrix}
1\\
-1\\
0
\end{pmatrix}.
\]

Therefore,

\[
\begin{aligned}
7a+7b+3c&=1,\\
-9a-9b-4c&=-1,\\
4a+4b+2c&=0.
\end{aligned}
\]

Again let

\[
s=a+b.
\]

Then

\[
\begin{aligned}
7s+3c&=1,\\
-9s-4c&=-1,\\
4s+2c&=0.
\end{aligned}
\]

From the third equation,

\[
4s+2c=0.
\]

Therefore,

\[
c=-2s.
\]

Substitute into the first equation:

\[
7s+3(-2s)=1.
\]

Hence

\[
7s-6s=1,
\]

so

\[
s=1.
\]

Therefore,

\[
c=-2.
\]

Since

\[
a+b=1,
\]

we are free to choose any convenient pair satisfying this relation.

Choose

\[
a=1,
\qquad
b=0.
\]

Then

\[
\boxed{
v_2=
\begin{pmatrix}
1\\
0\\
-2
\end{pmatrix}.
}
\]

Let us immediately verify this choice:

\[
Nv_2
=
\begin{pmatrix}
7&7&3\\
-9&-9&-4\\
4&4&2
\end{pmatrix}
\begin{pmatrix}
1\\
0\\
-2
\end{pmatrix}.
\]

Carrying out the multiplication,

\[
Nv_2
=
\begin{pmatrix}
7(1)+7(0)+3(-2)\\
-9(1)-9(0)-4(-2)\\
4(1)+4(0)+2(-2)
\end{pmatrix}.
\]

Thus

\[
Nv_2
=
\begin{pmatrix}
7-6\\
-9+8\\
4-4
\end{pmatrix}
=
\begin{pmatrix}
1\\
-1\\
0
\end{pmatrix}.
\]

Therefore,

\[
\boxed{
Nv_2=v_1.
}
\]


% -------------------------------------------------
\subsection{Step 4: Search for the Parent of \(v_2\)}
% -------------------------------------------------

We still need a third independent vector.

Following exactly the same idea, we now ask whether \(v_2\) itself has
a parent.

We seek \(v_3\) satisfying

\[
\boxed{
Nv_3=v_2.
}
\]

Let

\[
v_3=
\begin{pmatrix}
a\\
b\\
c
\end{pmatrix}.
\]

Then

\[
\begin{pmatrix}
7&7&3\\
-9&-9&-4\\
4&4&2
\end{pmatrix}
\begin{pmatrix}
a\\
b\\
c
\end{pmatrix}
=
\begin{pmatrix}
1\\
0\\
-2
\end{pmatrix}.
\]

Therefore,

\[
\begin{aligned}
7a+7b+3c&=1,\\
-9a-9b-4c&=0,\\
4a+4b+2c&=-2.
\end{aligned}
\]

Again define

\[
s=a+b.
\]

Then

\[
\begin{aligned}
7s+3c&=1,\\
-9s-4c&=0,\\
4s+2c&=-2.
\end{aligned}
\]

From the second equation,

\[
-9s-4c=0.
\]

Therefore,

\[
4c=-9s,
\]

so

\[
c=-\frac{9}{4}s.
\]

Now substitute this into the third equation:

\[
4s+2\left(-\frac94s\right)=-2.
\]

Thus

\[
4s-\frac92s=-2.
\]

Writing \(4s\) as \(\frac82s\),

\[
\frac82s-\frac92s=-2.
\]

Therefore,

\[
-\frac12s=-2.
\]

Multiplying both sides by \(-2\),

\[
s=4.
\]

Hence

\[
c=-\frac94(4),
\]

so

\[
c=-9.
\]

Since

\[
a+b=4,
\]

choose

\[
a=4,
\qquad
b=0.
\]

Therefore,

\[
\boxed{
v_3=
\begin{pmatrix}
4\\
0\\
-9
\end{pmatrix}.
}
\]

Again we verify directly:

\[
Nv_3
=
\begin{pmatrix}
7&7&3\\
-9&-9&-4\\
4&4&2
\end{pmatrix}
\begin{pmatrix}
4\\
0\\
-9
\end{pmatrix}.
\]

Therefore,

\[
Nv_3
=
\begin{pmatrix}
7(4)+7(0)+3(-9)\\
-9(4)-9(0)-4(-9)\\
4(4)+4(0)+2(-9)
\end{pmatrix}.
\]

Hence

\[
Nv_3
=
\begin{pmatrix}
28-27\\
-36+36\\
16-18
\end{pmatrix}
=
\begin{pmatrix}
1\\
0\\
-2
\end{pmatrix}.
\]

Thus

\[
\boxed{
Nv_3=v_2.
}
\]


% -------------------------------------------------
\subsection{Step 5: We Have Found the Jordan Chain}
% -------------------------------------------------

We have now obtained

\[
v_1=
\begin{pmatrix}
1\\
-1\\
0
\end{pmatrix},
\qquad
v_2=
\begin{pmatrix}
1\\
0\\
-2
\end{pmatrix},
\qquad
v_3=
\begin{pmatrix}
4\\
0\\
-9
\end{pmatrix}.
\]

They satisfy

\[
Nv_1=0,
\]

\[
Nv_2=v_1,
\]

and

\[
Nv_3=v_2.
\]

Since

\[
N=A-2I,
\]

the chain is

\[
\boxed{
v_3
\xrightarrow{\,A-2I\,}
v_2
\xrightarrow{\,A-2I\,}
v_1
\xrightarrow{\,A-2I\,}
0.
}
\]

This is a Jordan chain of length \(3\).


% -------------------------------------------------
\subsection{Step 6: Recover the Action of \(A\)}
% -------------------------------------------------

From

\[
(A-2I)v_1=0,
\]

we have

\[
Av_1-2Iv_1=0.
\]

Since

\[
Iv_1=v_1,
\]

we obtain

\[
Av_1-2v_1=0,
\]

and hence

\[
\boxed{
Av_1=2v_1.
}
\]

Next,

\[
(A-2I)v_2=v_1.
\]

Therefore,

\[
Av_2-2Iv_2=v_1.
\]

Since

\[
Iv_2=v_2,
\]

we obtain

\[
Av_2-2v_2=v_1.
\]

Hence

\[
\boxed{
Av_2=2v_2+v_1.
}
\]

Finally,

\[
(A-2I)v_3=v_2.
\]

Therefore,

\[
Av_3-2Iv_3=v_2.
\]

Since

\[
Iv_3=v_3,
\]

we obtain

\[
Av_3-2v_3=v_2.
\]

Thus

\[
\boxed{
Av_3=2v_3+v_2.
}
\]

Altogether,

\[
\boxed{
\begin{aligned}
Av_1&=2v_1,\\
Av_2&=v_1+2v_2,\\
Av_3&=v_2+2v_3.
\end{aligned}
}
\]


% -------------------------------------------------
\subsection{Step 7: Construct the Jordan Matrix}
% -------------------------------------------------

The preceding relations immediately tell us what the matrix of \(A\)
must look like in the ordered basis

\[
\mathcal B=(v_1,v_2,v_3).
\]

From

\[
Av_1=2v_1,
\]

the first coordinate column is

\[
\begin{pmatrix}
2\\
0\\
0
\end{pmatrix}.
\]

From

\[
Av_2=v_1+2v_2,
\]

the second coordinate column is

\[
\begin{pmatrix}
1\\
2\\
0
\end{pmatrix}.
\]

From

\[
Av_3=v_2+2v_3,
\]

the third coordinate column is

\[
\begin{pmatrix}
0\\
1\\
2
\end{pmatrix}.
\]

Therefore,

\[
\boxed{
J=
\begin{pmatrix}
2&1&0\\
0&2&1\\
0&0&2
\end{pmatrix}.
}
\]


% -------------------------------------------------
\subsection{Step 8: Construct the Matrix \(P\)}
% -------------------------------------------------

The columns of \(P\) are the vectors of the Jordan chain in the order

\[
v_1,\;v_2,\;v_3.
\]

Therefore,

\[
P=
\begin{pmatrix}
\vert&\vert&\vert\\
v_1&v_2&v_3\\
\vert&\vert&\vert
\end{pmatrix}.
\]

Substituting the vectors,

\[
\boxed{
P=
\begin{pmatrix}
1&1&4\\
-1&0&0\\
0&-2&-9
\end{pmatrix}.
}
\]

Before using \(P^{-1}\), we verify that \(P\) is invertible.

Compute

\[
\det P
=
\det
\begin{pmatrix}
1&1&4\\
-1&0&0\\
0&-2&-9
\end{pmatrix}.
\]

Expanding along the second row is convenient:

\[
\det P
=
(-1)(-1)^{2+1}
\begin{vmatrix}
1&4\\
-2&-9
\end{vmatrix}.
\]

Thus

\[
\det P
=
\begin{vmatrix}
1&4\\
-2&-9
\end{vmatrix}.
\]

Therefore,

\[
\det P
=
1(-9)-4(-2).
\]

Hence

\[
\det P=-9+8=-1.
\]

Thus

\[
\boxed{
\det P=-1\neq0.
}
\]

Therefore \(P\) is invertible.


% -------------------------------------------------
\subsection{Step 9: Verify \(AP=PJ\)}
% -------------------------------------------------

Before computing \(P^{-1}AP\), it is useful to verify the equivalent
relation

\[
AP=PJ.
\]

First compute \(AP\):

\[
AP
=
\begin{pmatrix}
9&7&3\\
-9&-7&-4\\
4&4&4
\end{pmatrix}
\begin{pmatrix}
1&1&4\\
-1&0&0\\
0&-2&-9
\end{pmatrix}.
\]

Compute the entries row by row.

For the first row,

\[
\begin{aligned}
(1,1)&:\quad 9(1)+7(-1)+3(0)=2,\\
(1,2)&:\quad 9(1)+7(0)+3(-2)=3,\\
(1,3)&:\quad 9(4)+7(0)+3(-9)=9.
\end{aligned}
\]

For the second row,

\[
\begin{aligned}
(2,1)&:\quad -9(1)-7(-1)-4(0)=-2,\\
(2,2)&:\quad -9(1)-7(0)-4(-2)=-1,\\
(2,3)&:\quad -9(4)-7(0)-4(-9)=0.
\end{aligned}
\]

For the third row,

\[
\begin{aligned}
(3,1)&:\quad 4(1)+4(-1)+4(0)=0,\\
(3,2)&:\quad 4(1)+4(0)+4(-2)=-4,\\
(3,3)&:\quad 4(4)+4(0)+4(-9)=-20.
\end{aligned}
\]

Therefore,

\[
\boxed{
AP=
\begin{pmatrix}
2&3&9\\
-2&-1&0\\
0&-4&-20
\end{pmatrix}.
}
\]

Now compute \(PJ\):

\[
PJ
=
\begin{pmatrix}
1&1&4\\
-1&0&0\\
0&-2&-9
\end{pmatrix}
\begin{pmatrix}
2&1&0\\
0&2&1\\
0&0&2
\end{pmatrix}.
\]

Again compute every entry.

For the first row,

\[
\begin{aligned}
(1,1)&:\quad 1(2)+1(0)+4(0)=2,\\
(1,2)&:\quad 1(1)+1(2)+4(0)=3,\\
(1,3)&:\quad 1(0)+1(1)+4(2)=9.
\end{aligned}
\]

For the second row,

\[
\begin{aligned}
(2,1)&:\quad (-1)(2)+0(0)+0(0)=-2,\\
(2,2)&:\quad (-1)(1)+0(2)+0(0)=-1,\\
(2,3)&:\quad (-1)(0)+0(1)+0(2)=0.
\end{aligned}
\]

For the third row,

\[
\begin{aligned}
(3,1)&:\quad 0(2)+(-2)(0)+(-9)(0)=0,\\
(3,2)&:\quad 0(1)+(-2)(2)+(-9)(0)=-4,\\
(3,3)&:\quad 0(0)+(-2)(1)+(-9)(2)=-20.
\end{aligned}
\]

Therefore,

\[
\boxed{
PJ=
\begin{pmatrix}
2&3&9\\
-2&-1&0\\
0&-4&-20
\end{pmatrix}.
}
\]

Hence

\[
\boxed{
AP=PJ.
}
\]

This already verifies that the Jordan-chain construction is correct.


% -------------------------------------------------
\subsection{Step 10: Verify \(P^{-1}AP=J\)}
% -------------------------------------------------

For completeness, we now verify the similarity transformation directly.

For this matrix,

\[
P^{-1}
=
\begin{pmatrix}
0&-1&0\\
9&9&4\\
-2&-2&-1
\end{pmatrix}.
\]

We already found

\[
AP=
\begin{pmatrix}
2&3&9\\
-2&-1&0\\
0&-4&-20
\end{pmatrix}.
\]

Therefore,

\[
P^{-1}AP
=
\begin{pmatrix}
0&-1&0\\
9&9&4\\
-2&-2&-1
\end{pmatrix}
\begin{pmatrix}
2&3&9\\
-2&-1&0\\
0&-4&-20
\end{pmatrix}.
\]

Compute the first row:

\[
\begin{aligned}
(1,1)&:\quad 0(2)+(-1)(-2)+0(0)=2,\\
(1,2)&:\quad 0(3)+(-1)(-1)+0(-4)=1,\\
(1,3)&:\quad 0(9)+(-1)(0)+0(-20)=0.
\end{aligned}
\]

Compute the second row:

\[
\begin{aligned}
(2,1)&:\quad 9(2)+9(-2)+4(0)=0,\\
(2,2)&:\quad 9(3)+9(-1)+4(-4)=2,\\
(2,3)&:\quad 9(9)+9(0)+4(-20)=1.
\end{aligned}
\]

Compute the third row:

\[
\begin{aligned}
(3,1)&:\quad (-2)(2)+(-2)(-2)+(-1)(0)=0,\\
(3,2)&:\quad (-2)(3)+(-2)(-1)+(-1)(-4)=0,\\
(3,3)&:\quad (-2)(9)+(-2)(0)+(-1)(-20)=2.
\end{aligned}
\]

Therefore,

\[
P^{-1}AP
=
\begin{pmatrix}
2&1&0\\
0&2&1\\
0&0&2
\end{pmatrix}.
\]

Hence,

\[
\boxed{
P^{-1}AP
=
J
=
\begin{pmatrix}
2&1&0\\
0&2&1\\
0&0&2
\end{pmatrix}.
}
\]


% -------------------------------------------------
\subsection{What This Example Reveals}
% -------------------------------------------------

This single \(3\times3\) problem contains nearly the entire basic story
of the Jordan canonical form.

The characteristic polynomial told us that the eigenvalue \(2\)
appears three times:

\[
p_A(x)=(x-2)^3.
\]

However, the ordinary eigenspace supplied only one independent
eigenvector:

\[
v_1.
\]

Ordinary diagonalization therefore failed.

Rather than abandoning the problem, we traced backward:

\[
Nv_2=v_1,
\]

and then

\[
Nv_3=v_2.
\]

Thus the missing directions were recovered through the Jordan chain

\[
\boxed{
v_3
\longrightarrow
v_2
\longrightarrow
v_1
\longrightarrow
0.
}
\]

In the language of our earlier intuition, \(v_1\) was only the final
visible member of a larger mathematical family.  Searching for its
parent and grandparent recovered the vectors that ordinary
diagonalization could not provide.

Those vectors then formed the basis

\[
P=[\,v_1\;v_2\;v_3\,],
\]

and in that basis the complicated matrix

\[
A=
\begin{pmatrix}
9&7&3\\
-9&-7&-4\\
4&4&4
\end{pmatrix}
\]

became the remarkably simple matrix

\[
\boxed{
J=
\begin{pmatrix}
2&1&0\\
0&2&1\\
0&0&2
\end{pmatrix}.
}
\]

The diagonal entries preserve the eigenvalue \(2\), while the
superdiagonal \(1\)'s preserve the connections between successive
members of the Jordan chain.

Thus Jordan canonical form does not somehow manufacture the missing
eigenvectors.  Instead, it replaces the unavailable ordinary
eigendirections with generalized eigenvectors that retain the
remaining algebraic structure of the transformation.

% =================================================
% =================================================
\section{Computing the Jordan Canonical Form in MATLAB}
% =================================================

After working through the Jordan construction by hand, it is useful to
see how the same computation can be performed in MATLAB.

MATLAB provides the function

\begin{verbatim}
J = jordan(A)
\end{verbatim}

for computing the Jordan canonical form.  If we also want the
transformation matrix, we use

\begin{verbatim}
[P,J] = jordan(A)
\end{verbatim}

where the output satisfies

\[
\boxed{
P^{-1}AP=J.
}
\]

The columns of \(P\) form the corresponding Jordan basis, consisting
of eigenvectors and, when necessary, generalized eigenvectors.


% -------------------------------------------------
\subsection{Our \(3\times3\) Example in MATLAB}
% -------------------------------------------------

Consider again

\[
A=
\begin{pmatrix}
9&7&3\\
-9&-7&-4\\
4&4&4
\end{pmatrix}.
\]

The basic MATLAB calculation is remarkably short:

\begin{verbatim}
clear; clc;

A = [ 9  7  3;
     -9 -7 -4;
      4  4  4];

[P,J] = jordan(A);

disp('Transformation matrix P:')
disp(P)

disp('Jordan canonical form J:')
disp(J)
\end{verbatim}

For this problem, the Jordan matrix has the structure

\[
\boxed{
J=
\begin{pmatrix}
2&1&0\\
0&2&1\\
0&0&2
\end{pmatrix}.
}
\]

The matrix \(P\) returned by MATLAB need not contain exactly the same
generalized eigenvectors that we constructed by hand.  Generalized
eigenvectors are not unique.  What matters is that the returned
matrices satisfy

\[
\boxed{
P^{-1}AP=J.
}
\]

A convenient verification is obtained from the equivalent relation

\[
AP=PJ.
\]

Thus we may check

\begin{verbatim}
A*P - P*J
\end{verbatim}

and verify that the result is the zero matrix.

We may also directly examine

\begin{verbatim}
inv(P)*A*P
\end{verbatim}

and compare the result with \(J\).

For learning purposes, this is essentially all that is required:
construct \(A\), call \texttt{jordan}, examine \(P\) and \(J\), and
verify the decomposition.


% -------------------------------------------------
% -------------------------------------------------
% -------------------------------------------------
\subsection{Numerical Consideration}
% -------------------------------------------------

The preceding discussion concerns the exact algebraic structure of a
matrix.  In a textbook problem, the entries of \(A\) are assumed to be
known exactly, and questions such as algebraic multiplicity, geometric
multiplicity, and Jordan structure are therefore perfectly well
defined.

Matrices arising in scientific computation, however, are often less
idealized.  They may come from experimental measurements, parameter
estimation, numerical simulation, discretization of a differential
equation or partial differential equation, or previous floating-point calculations.  Consequently, the
matrix actually stored in a computer may differ slightly from the ideal
mathematical matrix.

A convenient way to represent this is

\[
\boxed{
A_{\varepsilon}=A+\varepsilon E,
}
\]

where \(\varepsilon\) is small and \(E\) describes the direction of the
perturbation.

The term

\[
\varepsilon E
\]

may be interpreted as a small amount of unavoidable impurity introduced
by measurement, modeling, discretization, or finite-precision
computation.


% -------------------------------------------------
\subsubsection*{Why This Matters for Jordan Form}
% -------------------------------------------------

Jordan structure is extremely sensitive to perturbations.

Suppose the exact matrix \(A\) has a repeated eigenvalue

\[
2,\;2,\;2
\]

and contains a nontrivial Jordan block.

A nearby matrix

\[
A_{\varepsilon}=A+\varepsilon E
\]

may instead have three slightly different eigenvalues:

\[
\boxed{
2,\;2,\;2
\quad
\longrightarrow
\quad
2+\delta_1,\;
2+\delta_2,\;
2+\delta_3,
}
\]

where the quantities \(\delta_1,\delta_2,\delta_3\) are small but
generally unequal.

If the perturbed eigenvalues become distinct, the nearby matrix becomes
diagonalizable.

Thus a matrix that is exactly defective in theory may appear
nondefective when represented by real data or floating-point numbers.

This is not a contradiction.  The two matrices are simply not
algebraically identical.


% -------------------------------------------------
\subsubsection*{Impurity in Real Data}
% -------------------------------------------------

Small perturbations may arise naturally from

\begin{itemize}
    \item measurement uncertainty,
    \item sensor noise,
    \item parameter estimation,
    \item floating-point representation,
    \item rounding error,
    \item discretization error,
    \item or previous numerical computation.
\end{itemize}

Therefore, even if the ideal mathematical model contains an exact
Jordan block, the matrix available in a computer may instead be a
nearby matrix having a complete set of ordinary eigenvectors.

In some numerical calculations, this may make the nearby problem easier
to manipulate.

The important distinction is

\[
\boxed{
\text{numerically easier}
\neq
\text{algebraically identical}.
}
\]

The nearby matrix may be simpler computationally while no longer
possessing the exact Jordan structure of the original matrix.


% -------------------------------------------------
\subsection{A Simple MATLAB Experiment}
% -------------------------------------------------

This sensitivity can be demonstrated directly in MATLAB.

Consider again

\begin{verbatim}
clear; clc;

A = [ 9  7  3;
     -9 -7 -4;
      4  4  4];
\end{verbatim}

The exact matrix has the repeated eigenvalue \(2\).

Now suppose a small perturbation is already present:

\begin{verbatim}
rng(1);

epsilon = 1e-8;
E = randn(size(A));

Apert = A + epsilon*E;

disp('Eigenvalues of the original matrix:')
disp(eig(A))

disp('Eigenvalues of the perturbed matrix:')
disp(eig(Apert))
\end{verbatim}

The eigenvalues of the perturbed matrix will generally no longer be
exactly identical.

Even though

\[
\|A_{\varepsilon}-A\|
\]

is small, the algebraic classification of the matrix may change.

This illustrates one of the central numerical difficulties associated
with Jordan form:

\[
\boxed{
\text{small changes in the matrix can produce large changes in its
Jordan structure}.
}
\]


% -------------------------------------------------
\subsection{Exact Algebra Versus Floating-Point Computation}
% -------------------------------------------------

The Jordan canonical form therefore provides an excellent example of
the distinction between exact linear algebra and numerical linear
algebra.

In exact mathematics, one asks

\[
\boxed{
\text{What is the Jordan structure of }A?
}
\]

In numerical computation, one must also ask

\[
\boxed{
\text{How reliably can that structure be detected from the available
data?}
}
\]

A nontrivial Jordan block represents an exact algebraic relationship.
Because such relationships can be destroyed by very small perturbations,
they are difficult to identify reliably from approximate floating-point
data.

This is one reason that Jordan canonical form is extremely important
theoretically but is used with caution in numerical computation.


% -------------------------------------------------
\subsection{A Changing Computational Landscape}
% -------------------------------------------------

There is, however, an important broader perspective.

For several decades, mainstream scientific computing was dominated by
floating-point numerical methods, while exact symbolic manipulation
played a more specialized role.

In recent years, symbolic computation has gained renewed momentum.
Modern computer algebra systems, faster hardware, automatic algebraic
manipulation, exact arithmetic, and hybrid symbolic--numerical methods
have made exact computation practical in a much wider range of
problems.

This development is particularly relevant to questions such as Jordan
form, because the major difficulty often arises not from the algebra
itself, but from representing exact algebraic structure using
approximate numbers.

As symbolic and symbolic--numerical methods continue to improve, some
of the computational difficulties currently associated with detecting
exact Jordan structure may become less severe.

It is therefore reasonable to expect that future computational tools
may handle certain defective-matrix problems more naturally than
present floating-point methods do.

For now, however, the essential lesson is

\[
\boxed{
\begin{array}{c}
\text{exact algebra reveals the Jordan structure,}\\[1mm]
\text{floating-point perturbations may obscure or destroy it,}\\[1mm]
\text{and symbolic computation offers a possible bridge between the two.}
\end{array}
}
\]

The Jordan canonical form is therefore not only a classical theoretical
construction.  It is also a useful example of how the meaning of a
mathematical object can depend strongly on the computational framework
used to represent it.

% =================================================
\begin{comment}
\section{Conclusion}
% =================================================

We began this note with one of the most familiar ideas in linear
algebra: diagonalization.  When a matrix possesses enough linearly
independent eigenvectors, the situation is remarkably clean.  We form
a basis from those eigenvectors and obtain

\[
P^{-1}AP=D,
\]

where \(D\) is diagonal.

The difficulty appears when the eigenvectors run out.

A repeated eigenvalue may have algebraic multiplicity larger than the
dimension of its eigenspace, leaving us without enough ordinary
eigenvectors to construct a basis.  At first sight, this appears to be
the end of the diagonalization procedure.

But mathematics rarely progresses by simply declaring defeat when the
most obvious method stops working.

Instead, they changed the question.

Rather than asking only

\[
\boxed{
\text{Where does a vector go under }A-\lambda I?
}
\]

we also began asking

\[
\boxed{
\text{Where could this vector have come from?}
}
\]

That small change in perspective opened the door to the entire Jordan
construction.

Starting from an ordinary eigenvector

\[
(A-\lambda I)v_1=0,
\]

we searched backward for a vector \(v_2\) satisfying

\[
(A-\lambda I)v_2=v_1.
\]

Then we asked whether \(v_2\) itself had a predecessor:

\[
(A-\lambda I)v_3=v_2.
\]

Continuing in this manner produced the chain

\[
\boxed{
v_k
\longrightarrow
v_{k-1}
\longrightarrow
\cdots
\longrightarrow
v_2
\longrightarrow
v_1
\longrightarrow
0.
}
\]

What initially appeared to be a shortage of eigenvectors was therefore
not the absence of structure.  The structure was simply deeper than
ordinary eigenvectors alone could reveal.

This is one of the broader lessons of the Jordan canonical form.


% -------------------------------------------------
\subsection*{When the Forward Search Fails, Look Backward}
% -------------------------------------------------

Much of elementary mathematics trains us to think forward:

\[
\text{given something}
\quad\longrightarrow\quad
\text{compute what comes next}.
\]

That is an enormously useful habit, but it is not the only way
mathematicians think.

When a forward search becomes blocked, it can be fruitful to reverse
the direction of the question.

Instead of asking only

\[
\text{What does this object produce?}
\]

one may ask

\[
\text{What could have produced this object?}
\]

The Jordan chain is a particularly elegant example of this principle.

The ordinary eigenvector \(v_1\) appears to be the end of the story
because

\[
(A-\lambda I)v_1=0.
\]

But by searching backward, we discover that \(v_1\) may itself be the
image of \(v_2\), which may be the image of \(v_3\), and so forth.

The apparent endpoint becomes the starting point of a new
investigation.

This style of reasoning extends far beyond Jordan canonical form.
Mathematical problem solving often involves changing direction,
reversing a map, reconstructing causes from consequences, or asking
what hidden structure must exist in order to produce what we already
observe.


% -------------------------------------------------
\subsection*{Broadening Intuition}
% -------------------------------------------------

There is another lesson worth carrying beyond this particular topic.

At some point in mathematical study, intuition must become broader
than formal manipulation alone.

Proofs are indispensable.  Computation is indispensable.  Precise
definitions are indispensable.  But when a subject becomes buried
under increasingly abstruse notation or lengthy algebra, repeatedly
pushing symbols around may not be the best way to gain understanding.

Sometimes one should temporarily put the formalism aside and ask:

\begin{center}
\emph{What is actually happening here?}
\end{center}

That was the purpose of the practical analogies used in this note.
They were not substitutes for proof.  Rather, they provided another
way of seeing the structure that the proof was describing.

Once the underlying picture becomes visible, one can return to the
formal mathematics with a much better sense of why each equation is
there.

A useful habit is therefore

\[
\boxed{
\text{formal theory}
\;\longrightarrow\;
\text{intuition}
\;\longrightarrow\;
\text{return to the formal theory}.
}
\]

The second encounter with the mathematics is often much clearer than
the first.


% -------------------------------------------------
\end{comment}
% =================================================
\section{Conclusion}
% =================================================

We began this note with one of the most familiar ideas in linear
algebra: diagonalization.  When a matrix possesses enough linearly
independent eigenvectors, the situation is remarkably clean.  We form
a basis from those eigenvectors and obtain

\[
P^{-1}AP=D,
\]

where \(D\) is diagonal.

The difficulty appears when the eigenvectors run out.

A repeated eigenvalue may have algebraic multiplicity larger than the
dimension of its eigenspace, leaving us without enough ordinary
eigenvectors to construct a basis.  At first sight, this appears to be
the end of the diagonalization procedure.

But mathematics rarely progresses by simply declaring defeat when the
most obvious method stops working.

Instead, we changed the question.

Rather than asking only

\[
\boxed{
\text{Where does a vector go under }A-\lambda I?
}
\]

we also began asking

\[
\boxed{
\text{Where could this vector have come from?}
}
\]

That small change in perspective opened the door to the entire Jordan
construction.

Starting from an ordinary eigenvector

\[
(A-\lambda I)v_1=0,
\]

we searched backward for a vector \(v_2\) satisfying

\[
(A-\lambda I)v_2=v_1.
\]

Then we asked whether \(v_2\) itself had a predecessor:

\[
(A-\lambda I)v_3=v_2.
\]

Continuing in this manner produced the chain

\[
\boxed{
v_k
\longrightarrow
v_{k-1}
\longrightarrow
\cdots
\longrightarrow
v_2
\longrightarrow
v_1
\longrightarrow
0.
}
\]

What initially appeared to be a shortage of eigenvectors was therefore
not the absence of structure.  The structure was simply deeper than
ordinary eigenvectors alone could reveal.

This is one of the broader lessons of the Jordan canonical form.


% -------------------------------------------------
\subsection*{When the Forward Search Fails, Look Backward}
% -------------------------------------------------

Much of elementary mathematics trains us to think forward:

\[
\text{given something}
\quad\longrightarrow\quad
\text{compute what comes next}.
\]

That is an enormously useful habit, but it is not the only way
mathematicians think.

When a forward search becomes blocked, it can be fruitful to reverse
the direction of the question.

Instead of asking only

\[
\text{What does this object produce?}
\]

one may ask

\[
\text{What could have produced this object?}
\]

The Jordan chain is a particularly elegant example of this principle.

The ordinary eigenvector \(v_1\) appears to be the end of the story
because

\[
(A-\lambda I)v_1=0.
\]

But by searching backward, we ask whether \(v_1\) might itself be the
image of some \(v_2\), whether \(v_2\) might be the image of some
\(v_3\), and so forth.

The apparent endpoint becomes the starting point of a new
investigation.

This style of reasoning extends far beyond Jordan canonical form.
Mathematical problem solving often involves changing direction,
reversing a map, reconstructing causes from consequences, or asking
what hidden structure must exist in order to produce what we already
observe.


% -------------------------------------------------
\subsection*{Broadening Intuition}
% -------------------------------------------------

There is another lesson worth carrying beyond this particular topic.

At some point in mathematical study, intuition must become broader
than formal manipulation alone.

Proofs are indispensable.  Computation is indispensable.  Precise
definitions are indispensable.  But when a subject becomes buried
under increasingly abstruse notation or lengthy algebra, repeatedly
pushing symbols around may not be the best way to gain understanding.

Sometimes one should temporarily put the formalism aside and ask:

\begin{center}
\emph{What is actually happening here?}
\end{center}

That was the purpose of the practical analogies used in this note.
They were not substitutes for proof.  Rather, they provided another
way of seeing the structure that the proof was describing.

Once the underlying picture becomes visible, one can return to the
formal mathematics with a much better sense of why each equation is
there.

A useful habit is therefore

\[
\boxed{
\text{formal theory}
\;\longrightarrow\;
\text{intuition}
\;\longrightarrow\;
\text{return to the formal theory}.
}
\]

The second encounter with the mathematics is often much clearer than
the first.

The Jordan canonical form therefore leaves us with a lesson that is
larger than the decomposition itself.  When a familiar method reaches
an obstruction, the answer is not always to push the same calculation
harder.  Sometimes progress comes from changing the question, reversing
the direction of thought, or enlarging the framework in which the
problem is being viewed.

In our case, the failure of ordinary diagonalization was not the end
of the story.  It was precisely the obstruction that led us to
generalized eigenvectors, Jordan chains, and ultimately a richer
understanding of the transformation.

\begin{comment}
    

\addcontentsline{toc}{section}{References}

\begin{thebibliography}{9}

\bibitem{Finkbeiner}
Daniel T. Finkbeiner II,
\emph{Introduction to Matrices and Linear Transformations},
3rd ed.,
W. H. Freeman and Company,
San Francisco.
\end{thebibliography}
\end{comment}
% =================================================
\addcontentsline{toc}{section}{References}

\begin{thebibliography}{9}

\bibitem{Finkbeiner}
Daniel T. Finkbeiner II,
\emph{Introduction to Matrices and Linear Transformations},
3rd ed.,
W. H. Freeman and Company,
San Francisco.

\bibitem{HornJohnson}
Roger A. Horn and Charles R. Johnson,
\emph{Matrix Analysis},
2nd ed.,
Cambridge University Press,
Cambridge,
2013.

\end{thebibliography}
\end{document}